\documentclass[11pt,a4paper]{article}
\usepackage[T1]{fontenc}
\usepackage[utf8]{inputenc}
\usepackage[english]{babel}
\usepackage{lmodern,microtype}
\usepackage[a4paper,margin=2.35cm,headheight=15pt]{geometry}
\usepackage{enumitem,booktabs,longtable,array,tabularx}
\usepackage{xcolor,hyperref,xurl,fancyhdr,titlesec,listings,framed}
\definecolor{ddtadark}{RGB}{28,38,52}
\definecolor{ddtagray}{RGB}{242,244,247}
\definecolor{ddtablue}{RGB}{42,88,135}
\definecolor{ddtagreen}{RGB}{48,111,78}
\definecolor{ddtaorange}{RGB}{148,91,25}
\hypersetup{colorlinks=true,linkcolor=ddtablue,urlcolor=ddtablue,pdftitle={DDTA Research Work Plan after Documentation Closure R2},pdfauthor={Documentation-Driven Threat Analysis research workspace}}
\pagestyle{fancy}\fancyhf{}\lhead{DDTA research work plan}\rhead{Revision 2}\cfoot{\thepage}
\titleformat{\section}{\Large\bfseries\color{ddtadark}}{\thesection}{0.7em}{}
\titleformat{\subsection}{\large\bfseries\color{ddtadark}}{\thesubsection}{0.7em}{}
\setlist[itemize]{leftmargin=1.55em,itemsep=0.25em,topsep=0.35em}
\setlist[enumerate]{leftmargin=1.8em,itemsep=0.25em,topsep=0.35em}
\setlength{\parskip}{0.35em}\setlength{\parindent}{0pt}\setlength{\emergencystretch}{2em}
\lstset{basicstyle=\ttfamily\small,columns=fullflexible,keepspaces=true,frame=single,framerule=0.3pt,rulecolor=\color{ddtagray},backgroundcolor=\color{ddtagray},xleftmargin=0.8em,xrightmargin=0.8em,breaklines=true}
\newenvironment{gate}[1]{\def\FrameCommand{\fcolorbox{ddtagreen}{ddtagray}}\MakeFramed{\advance\hsize-2\width\FrameRestore}\noindent\textbf{#1}\par\smallskip}{\endMakeFramed}
\newenvironment{boundary}[1]{\def\FrameCommand{\fcolorbox{ddtaorange}{ddtagray}}\MakeFramed{\advance\hsize-2\width\FrameRestore}\noindent\textbf{#1}\par\smallskip}{\endMakeFramed}
\begin{document}
\begin{titlepage}\centering\vspace*{2cm}{\Large Documentation-Driven Threat Analysis\par}\vspace{0.8cm}{\Huge\bfseries Research Work Plan\par}\vspace{0.3cm}{\LARGE After documentation-layer closure -- Revision 2\par}\vspace{1cm}{\large From the frozen documentation metamodel to Base Analysis, BAE and multi-method analysis\par}\vfill\begin{minipage}{0.84\textwidth}\small
\textbf{Baseline.} \texttt{ec6c0107e0d5c0460d2afcbe12dae6f40bc5e6c5}.\\[0.5em]
\textbf{Status.} Current execution plan after the Chapters 2--4 documentation-layer freeze. W0 is CLOSED and BA0 is NEXT. It supersedes Revision 1 sequencing but does not reopen closed documentation semantics.\end{minipage}\vfill{\large 15 August 2026\par}\end{titlepage}
\tableofcontents\newpage
\section{Current checkpoint}
The governed documentation metamodel is closed for the current thesis scope through:
\begin{lstlisting}
Project problem framing [method precondition]
        -> MacroRequirement
        -> Decision
        -> FunctionalRequirement
        -> SpecializedRequirement [abstract]
             -> SecurityRequirement
\end{lstlisting}
The common requirement abstraction is also closed:
\begin{lstlisting}
GovernedDocument
        |
        `-- Requirement [abstract]
                normativeClause : NormativeClause [1..*]
                |
                +-- FunctionalRequirement
                `-- SpecializedRequirement [abstract]
                        `-- SecurityRequirement
\end{lstlisting}
\texttt{Requirement} is the governed normative obligation. The separate L1 \texttt{NormativeObligation} wrapper is rejected. SecurityRequirement adds \texttt{protectedSecurityProperty : SecurityProperty [1]}; its security failure mode must be identifiable in inherited normative clauses and its identity is cause-neutral.

\section{Thesis-writing freeze gate: Chapters 2--4}
\subsection{W0.1 -- Chapter 2 -- CLOSED}
Chapter 2 is semantically closed. Only typographical/citation correction, stale forward-reference cleanup and terminology normalization that does not change claims are allowed.
\subsection{W0.2 -- Chapter 3 -- CLOSED}
Chapter 3 is semantically closed. G1--G4 and the thesis research scope remain unchanged; finalization is editorial only.
\subsection{W0.3 -- Chapter 4 -- CLOSED}
The finalization commit consolidates the already closed results reached after S1R1: Requirement abstraction, normative clauses, rejection of NormativeObligation, S2 SecurityRequirement, protected property, failure-mode explicitness, cause neutrality, single governing security property, and negative boundaries.
\begin{gate}{W0 exit gate}
W0 is CLOSED: Chapters 2, 3 and 4 are CLOSED / FINAL for the current thesis scope and are not routinely edited during Base Analysis, Finding, overlay or STRIDE work. The next authorized microstep is BA0.
\end{gate}

\section{Reopen policy for frozen chapters}
\subsection{Chapter 2}
Reopen only for a material literature/factual error, a materially unsupported claim, or falsification of a foundational distinction used by the thesis.
\subsection{Chapter 3}
Reopen only if the research gap or evaluated thesis scope is materially falsified/changed, or prior work is found to solve a claimed discontinuity in the evaluated end-to-end sense.
\subsection{Chapter 4}
Reopen only for a concrete recurring downstream counterexample that cannot be represented without changing the governed documentation contract. A new BAE relation, plugin-private taxonomy field, vocabulary suggestion or implementation control is not by itself a reopen trigger.
\begin{boundary}{Reopen discipline}
Every reopening must name the counterexample, modify the minimum owning layer and regress the previously accepted corpora.
\end{boundary}

\section{Base Analysis block}
\subsection{BA0 -- Base Analysis responsibility and boundary}
Define Base Analysis as a governed, methodology-neutral analyzable representation of system/project knowledge grounded in governed documentation and explicit reviewed additions. Falsify these boundaries:
\begin{itemize}
\item not a second narrative documentation hierarchy;
\item not a STRIDE DFD by definition;
\item not automatically canonized NLP/LLM extraction;
\item may integrate reviewed structure derived from several governed documents;
\item one canonical knowledge base may support several projections;
\item an analysis method consumes the representation but cannot silently redefine its semantic core.
\end{itemize}
\begin{gate}{BA0 exit gate}A precise responsibility statement and non-goals are accepted before BAE taxonomy design.\end{gate}

\subsection{BA1 -- Minimal BAE ontology}
Start from the smallest vocabulary forced by the facial-access corpus:
\begin{lstlisting}
BAE [abstract]
    +-- Actor
    +-- Asset
    `-- Interaction
\end{lstlisting}
Pressure-test whether first-class Component/ServiceCapability, DataStore, CommunicationChannel, InterfaceContract, Boundary, Location/Environment or State/Condition are actually necessary. These are hypotheses, not pre-approved metaclasses.
\begin{gate}{BA1 exit gate}Minimal generic BAE set with definitions, stable identity rules and negative controls.\end{gate}

\subsection{BA2 -- Relations and canonical action vocabulary}
Define the smallest stable relation vocabulary needed for a graph without uncontrolled synonym growth. Candidate pressure points: source/target participation, produces/consumes/handles, stores/retrieves, dependsOn, boundary relations, invokes/delivers/receives and capability provision/consumption.
\begin{lstlisting}
L1 semantic relation
    !=
L2 canonical lexical label
    !=
L3 authoring suggestion / synonym normalization
\end{lstlisting}
Avoid a generic \texttt{affects}-style escape relation and avoid freezing a universal verb taxonomy prematurely.

\subsection{BA3 -- Document-to-BAE derivation, provenance and authority}
Required distinctions:
\begin{lstlisting}
governed document fact
candidate BAE extraction / suggestion
reviewed canonical BAE
reviewed analytical addition
\end{lstlisting}
Close stable identity, source references, authoritative/supporting sources, no automatic noun-to-BAE canonization, no automatic synonym merge, explicit candidate acceptance/rejection, BAE-to-document traceability, and explicit reviewed additions.

\subsection{BA4 -- Multi-level views and projections}
Views are projections of one canonical Base Analysis, not competing authored models. Candidate views:
\begin{enumerate}
\item Context / actor view.
\item Interaction view.
\item Asset / data-flow view.
\item Boundary / deployment view, only if boundary is canonical BAE semantics.
\item Security overlay view projecting SecurityRequirements/properties without turning Base Analysis into a threat-method model.
\end{enumerate}
A view may omit information for readability but must not imply that omitted information does not exist.

\subsection{BA5 -- Controlled vocabulary and authoring assistance}
Tooling may suggest reuse of Actor/Asset/BAE names, canonical actions/predicates, probable duplicates/paraphrases, taxonomy entries, BAE candidates/relations and existing references. Similarity or an LLM output never silently establishes identity, parentage, ownership or canonical BAE status.
\begin{lstlisting}
LLM-assisted authoring
    -> vocabulary / BAE suggestions
    -> reviewed governed documentation

LLM-assisted analysis preparation
    -> candidate BAE / relation extraction
    -> reviewed Base Analysis
\end{lstlisting}
\begin{gate}{BA5 exit gate}Assistance reduces synonym/paraphrase drift without making the tool semantic authority.\end{gate}

\subsection{BA6 -- Base Analysis closure and regression}
Regress closed documentation corpora and at least one structurally different holdout. Verify identity stability under paraphrase, mutation handling, view consistency, absence of method-specific threat categories in the core and explicit diagnostics for insufficient documentation.
\begin{gate}{Base Analysis exit gate}
A methodology-neutral Base Analysis can be built and reviewed from governed documentation plus explicit reviewed additions, with stable BAE identities, relations, provenance and projections.
\end{gate}

\section{Generic analysis envelope}
\subsection{A1 -- AnalysisRecord}Define common identity, method/version, scope, baseline and execution/provenance envelope with method-owned payload remaining plugin-owned.
\subsection{A2 -- Finding}Define a reviewed common result envelope that can reference affected BAE and governed requirements without copying method taxonomy into the common core.
\subsection{A3 -- Change/provenance integration}Resolve S1.5 many-to-many analysis/document history and distinguish analysis scope from accepted change. Test whether a general change-event abstraction is actually necessary.

\section{Method-neutral overlay and concrete evaluations}
\subsection{O1 -- Generic overlay contract}Define compatibility, scope selection, required BAE/view capabilities, method-owned payload, diagnostics, finding derivation and failure behavior.
\subsection{O2 -- STRIDE}Use STRIDE only after Base Analysis and O1 are stable enough for falsification.
\subsection{O3 -- Closed-loop design improvement}
\begin{lstlisting}
governed documentation
 -> Base Analysis / views
 -> analysis overlay
 -> AnalysisRecord
 -> reviewed Finding
 -> accepted governed change / SecurityRequirement
 -> Base Analysis refresh
 -> re-analysis
\end{lstlisting}
\subsection{O4 -- STRIDE-AI}Reuse the same common contracts without copying STRIDE-AI-private asset/failure/taxonomy semantics into the common core.

\section{Later neutrality challenges and ThreatForge}
Use LINDDUN and optionally another method as neutrality challenges after STRIDE/STRIDE-AI; do not claim universal neutrality from two implementations. ThreatForge remains an implementation/tooling case study and never establishes DDTA semantic authority.

\section{Next authorized microstep}
\begin{gate}{Authorized next step}
After the W0 finalization commit, execute only \textbf{BA0 -- Base Analysis responsibility and boundary}, immediately pressure-tested on the facial-access camera corpus. Do not start STRIDE, AnalysisRecord, Finding or broad BAE taxonomy enumeration before BA0 is closed.
\end{gate}
\end{document}
